
\subsection{School-Level Inflation with an Anchor Equation}
\label{sec:baseline_model}

Let $Y_{i,j,k,t}$ be a binary indicator that equals 1 if student $i$ in school $j$, taking subject $k$ at time $t$, achieves an $A$ or $A^*$, and 0 otherwise. We assume there is an unobserved latent variable $Y_{i,j,k,t}^*$ such that
\begin{equation}
    \label{eq:indicator}
    Y_{i,j,k,t}
    \;=\;
    \mathbf{1}\{Y_{i,j,k,t}^* \,\ge\, 0\}.
\end{equation}

\subsubsection*{Main Equation (Potentially Affected by Treatment)}

We specify the latent index for our \emph{primary} outcome ($Y$) as follows:
\begin{align}
    \label{eq:latent_index}
    Y_{i,j,k,t}^*
    \;=\;& \beta_0 
    \;+\;
    \beta_1\,\text{Year2020}_t
    \;+\;
    \sum_{j}\alpha_j\,\text{School}_j 
    \;+\;
    \sum_{j}\delta_j\,\bigl(\text{School}_j \times \text{Year2020}_t\bigr)
    \nonumber \\[-6pt]
    &\;+\;\sum_{k}\gamma_k\,\text{Subject}_k
    \;+\;
    \sum_{l}\eta_{l}\,\text{GCSE}_{i,l,t}
    \;+\;
    \sum_{m}\zeta_m\,X_{i,m}
    \;+\;
    \sum_{n}\nu_n\,X_{j,n}
    \;+\;
    \varepsilon_{i,j,k,t},
\end{align}
where:
\begin{itemize}
    \item \(\text{Year2020}_t = 1\) if $t=2020$ (pandemic year), and 0 otherwise,
    \item \(\text{School}_j\) is a set of school fixed effects (one for each school \(j\)),
    \item \(\text{School}_j \times \text{Year2020}_t\) allows each school’s effect to differ in 2020,
    \item \(\text{Subject}_k\) are subject fixed effects (one for each subject \(k\)), 
    \item \(\text{GCSE}_{i,l,t}\) are dummy variables for the student’s GCSE performance (e.g., by grade bins) in required subjects,
    \item \(X_{i,m}\) are additional student-level controls (e.g., gender, race, SES),
    \item \(X_{j,n}\) are additional school-level controls (e.g., school type, quality),
    \item \(\varepsilon_{i,j,k,t}\) is the error term, which we allow to be \emph{heteroskedastic} (see below).
\end{itemize}
Here, $\delta_j$ measures the *incremental* effect of being in school $j$ \emph{during} 2020 (relative to that school’s baseline).

\subsubsection*{Anchor Equation (Unaffected by Treatment)}

To aid identification of the GCSE slopes \(\eta_l\) (i.e., their absolute scale), we include a \emph{secondary} binary outcome $W_{i,j,t}$ that is ex-ante \emph{unaffected} by the treatment/pandemic grading policy. Denote the latent variable $W_{i,j,t}^*$, with
\[
    W_{i,j,t} \;=\; \mathbf{1}\{W_{i,j,t}^* \,\ge\, 0\}.
\]
We assume:
\begin{align}
\label{eq:anchor_index}
    W_{i,j,t}^*
    \;=\;
    \lambda_0
    \;+\;
    \sum_{l}\underbrace{\eta_{l}}_{\text{same GCSE slope}}\,\text{GCSE}_{i,l,t}
    \;+\;
    \sum_{r}\omega_r\,Z_{i,j,r}
    \;+\;
    \varepsilon_{i,j,t}^{(w)},
\end{align}
where
\begin{itemize}
    \item We \textbf{force} the GCSE coefficients in \eqref{eq:anchor_index} to be the \emph{same} \(\eta_l\) as in \eqref{eq:latent_index}, ensuring a common scale for these controls.
    \item $Z_{i,j,r}$ are additional covariates relevant for the anchor outcome but \emph{not} necessarily tied to the treatment.
    \item We assume $\varepsilon_{i,j,t}^{(w)} \sim N(0,\;\sigma_{w,j,t}^2)$; possibly heteroskedastic in a similar or simpler fashion.
    \item Critically, $W_{i,j,t}$ is not altered by $\text{Year2020}_t$ or $\delta_j$; that is, \(\text{Year2020}_t\) does \emph{not} enter \eqref{eq:anchor_index}, ensuring $W$ is an “anchor” unaffected by treatment.
\end{itemize}

\subsection{Estimating School-Specific Grade Inflation}

We continue to define *grade inflation* for school $j$ as the change in the predicted probability of an $A$ (or $A^*$) between the baseline (non-pandemic) period and the pandemic year 2020, holding student and school characteristics constant. Specifically, once we estimate all parameters in \eqref{eq:latent_index} and \eqref{eq:anchor_index} via maximum likelihood, we compute:
\begin{equation}
\label{eq:grade_inflation}
\theta_j \;=\;
\Pr\bigl(Y=1 \mid \text{Year2020}=1,\;\text{School}=j,\;\overline{X}\bigr)
\;-\;
\Pr\bigl(Y=1 \mid \text{Year2020}=0,\;\text{School}=j,\;\overline{X}\bigr),
\end{equation}
where $\overline{X}$ denotes a chosen reference vector of student and school characteristics. Since we use a probit link (see below), the predicted probability is
\[
\Pr(Y=1 \mid X) \;=\; \Phi\Bigl(\tfrac{X\hat{\beta}}{\sigma_i}\Bigr),
\]
where $\Phi(\cdot)$ is the standard normal CDF and $X\hat{\beta}/\sigma_i$ arises from \eqref{eq:latent_index}.

\subsection{Heteroskedastic Probit Estimation with a Joint Likelihood}

We estimate the parameters of both the \emph{main equation} \eqref{eq:latent_index} and the \emph{anchor equation} \eqref{eq:anchor_index} jointly via a \textbf{heteroskedastic probit} framework, allowing $\sigma_{j,t}$ to vary by school and year in the main equation, and similarly $\sigma_{w,j,t}$ for the anchor equation if needed:
\[
    \varepsilon_{i,j,k,t} = \sigma_{j,t}\,u_{i,j,k,t},\quad
    \varepsilon_{i,j,t}^{(w)} = \sigma_{w,j,t}\,u_{i,j,t}^{(w)},
    \quad
    u_{i,j,k,t},\,u_{i,j,t}^{(w)} \,\sim\, N(0,\,1).
\]
One might specify, for instance:
\[
    \sigma_{j,t} \;=\;\exp(Z_{j,t}\,\gamma),
    \quad
    \sigma_{w,j,t} \;=\;\exp(Z_{w,j,t}\,\gamma_w),
\]
with separate or overlapping covariates $Z_{j,t}$ and $Z_{w,j,t}$. 

\subsubsection*{Joint Log-Likelihood}
Define the observation-level contributions to each equation:
\[
\ell_Y(\beta,\gamma)
=
\sum_{i=1}^n
\Bigl[
Y_{i}\,\ln \Phi\!\Bigl(\tfrac{X_i\beta}{\sigma_i}\Bigr)
\;+\;
(1 - Y_{i}) \,\ln \Bigl(1 - \Phi\!\bigl(\tfrac{X_i\beta}{\sigma_i}\bigr)\Bigr)
\Bigr],
\]
\[
\ell_W(\eta,\omega,\gamma_w)
=
\sum_{i=1}^n
\Bigl[
W_{i}\,\ln \Phi\!\Bigl(\tfrac{X_i^{(w)}\tilde{\beta}}{\sigma_{w,i}}\Bigr)
\;+\;
(1 - W_{i}) \,\ln \Bigl(1 - \Phi\!\bigl(\tfrac{X_i^{(w)}\tilde{\beta}}{\sigma_{w,i}}\bigr)\Bigr)
\Bigr],
\]
where $X_i^{(w)}\tilde{\beta}$ includes the \(\eta_l\) terms (shared with the main equation’s GCSE slopes) plus additional $\omega_r$ parameters for $Z_{i,j,r}$. The total joint log-likelihood is:
\[
\mathcal{L}(\beta,\gamma,\eta,\omega,\gamma_w)
\;=\;
\ell_Y(\beta,\gamma)
\;+\;
\ell_W(\eta,\omega,\gamma_w),
\]
subject to the constraint that \(\eta_l\) in the anchor equation is the same as \(\eta_l\) in \eqref{eq:latent_index}. Maximizing $\mathcal{L}$ gives the parameter estimates $(\hat{\beta},\hat{\gamma},\hat{\eta},\hat{\omega},\hat{\gamma}_w)$, from which we compute school-specific inflation $\theta_j$ as in~\eqref{eq:grade_inflation}.

\bigskip
\noindent
\textbf{Remarks.}
\begin{itemize}
    \item This two-equation structure ensures that the \emph{GCSE slope parameters} $\eta_l$ are identified on a common scale between an \emph{unaffected} outcome (the anchor) and the \emph{treatment-affected} outcome ($Y$). 
    \item The heteroskedastic components $\sigma_{j,t}$ and $\sigma_{w,j,t}$ can thus be separated more cleanly from $\eta_l$, improving identification and interpretability of $\delta_j$ (the school-specific 2020 effect).
    \item If the anchor outcome $W_{i,j,t}$ is truly invariant to the 2020 treatment, it provides a stable “reference” to pin down how GCSE performance influences latent scores absent treatment-induced scaling changes.
\end{itemize}

\subsection{Simulation for school fixed effects}
